\documentclass[11pt,a4paper]{article}
\usepackage[utf8]{inputenc}
\usepackage[british]{babel}
\usepackage[a4paper,margin=2.5cm]{geometry}
\usepackage[hidelinks]{hyperref}
\setlength{\parindent}{0pt}\setlength{\parskip}{6pt}
\providecommand{\tightlist}{\setlength{\itemsep}{0pt}\setlength{\parskip}{0pt}}
\setcounter{secnumdepth}{0}
\title{\textit{The Ladies' Diary}\\[4pt]\large Nearly four hundred problems from an almanac for women, and the women who set them}
\author{Michael Leznik}\date{}
\begin{document}
\maketitle
In 1710 an almanac for women printed a puzzle written as a letter. A man
walking through Cheapside measures the shadow of Bow steeple at noon and
asks how tall it is. The first correct answer, printed the next year,
came from Mrs Mary Wright, who worked it out by trigonometry from the
height of the sun that day. The editor, John Tipper, sent her the twelve
free copies he had promised. Under her answer he admitted that he had
altered the length of the shadow on purpose, so that nobody could win by
looking the steeple's height up in a book. His readers were expected to
calculate. That was \emph{The Ladies' Diary}.

I decided to put my Claude account to work on something that might be
useful to everyone. After some thought, I chose \emph{The Ladies'
Diary}, an English almanack that ran from 1704 to 1840 and was famous
for its mathematical problems. The problems survive. You can find the
old volumes as scans online, and some OCR has been done on them. But in
practice they are almost unreadable. The scans are blurred, the OCR is
full of errors, the notation is two centuries out of date, and the
questions are written in the English of the 1700s, often in verse. The
mathematics is still good. It has just been locked away.

\hypertarget{a-magazine-for-women-that-printed-calculus}{%
\subsection{A magazine for women that printed
calculus}\label{a-magazine-for-women-that-printed-calculus}}

John Tipper, a schoolmaster in Coventry, launched the Diary in 1704. It
was a pocket almanack for women, with a calendar, recipes and useful
tables. It also had a page of riddles and mathematical questions,
answered by readers the following year, with the solvers' names printed
beside the best answers. Within a few years, the mathematics had taken
over.

Women set the questions as well as answering them. In 1714 a Mrs Barbara
Sidway asked for the largest cylinder that can be cut from a given cone.
The answer, which the Diary printed from four solvers, is to cut at one
third of the height, and the cylinder then holds exactly four ninths of
the cone. Today it is a standard exercise in calculus. In 1714 calculus
was still new. Lydia Fisher, Sarah Brown and Mary Nelson set questions
too.

Later came the men who became famous: Thomas Simpson, a weaver's son who
became mathematical master at the Royal Military Academy, Woolwich,
signed himself by comic pseudonyms, and edited the Diary from 1754 to
1760; William Emerson; John Landen, a land agent from Peterborough; and
Charles Hutton, who later edited the Diary himself. It asked for no
degree. It printed the best answer.

In 1817 Thomas Leybourn collected every mathematical question and answer
into four volumes. Those volumes are the main source for this project.
They stop in 1816, and for the last twenty-four years we went back to
the Diary itself.

\hypertarget{one-question-in-ten}{%
\subsection{One question in ten}\label{one-question-in-ten}}

Going through all four volumes, one question kept coming back. The Diary
was written for women. How many of the women actually did the
mathematics?

More than you might expect, and then far fewer.

Between 1709 and 1724 roughly one question in ten was set by a woman.
They answered, too. A small group appears again and again in those
years: Barbara Sidway, Mrs Mary Nelson of Newmarket, Mrs Anna Wright and
Mrs Mary Wright, and several who signed with pen-names: Silvia,
Hesychia, Anna Philomathes, Adrastea. The historian Shelley Costa spent
two years in English county record offices tracing some of these women.
That work is still unpublished.

Their questions are often about themselves. The Lady Mer. Heyshot asks
for the ages of her two sons: if the younger's age is taken from the
square of the elder's, 425 remains; the other way round, 235. (They are
21 and 16.) Mary Nelson describes a captain sharing prize money among
his crew, each man taking a pound more than the last and a hundredth
part of what remains, until it turns out every share is equal. (There
are 99 men, £99 each.)

Adrastea did something better. Her prize question for 1722 asks for the
depth of Eldon Hole, a pothole in the Peak District, from the time it
takes a dropped stone to be heard hitting the bottom. The printed answer
reports that on 18 August 1718 she really went there, dropped three
stones, and timed them: 9, 11 and 10 seconds. It also points out that
the method given in the standard reference book of the day was wrong. A
woman with a pendulum, standing over a hole in Derbyshire, correcting
the \emph{Lexicon Technicum}. That is an experiment.

Then it stops. After Adrastea's prize question of 1724, not one question
in Leybourn's first volume is set by a woman. The mathematics grows
harder, fluxions and curves and mechanics, and the names become those of
schoolmasters, surveyors and excisemen. Across the next century the
share of questions set by women runs at about one in sixty, then one in
eighty, then one in two hundred and twenty. The historian Shelley Costa
has described this as the masculinisation of the Diary's mathematical
forum. The numbers agree with her.

\hypertarget{the-ones-who-came-back}{%
\subsection{The ones who came back}\label{the-ones-who-came-back}}

They did not vanish entirely. In 1748 a writer signing herself Rosamond
asked for pairs of amicable numbers, and for the least number whose
divisors add up to seven more than itself. (It is 196.) In the same year
a Miss Munlove asked how many servants a gentleman keeps, if he can
place them in 960,799 different ways. The Diary printed the answer 7. It
was wrong. The formula it used counts arrangements with repetition,
which is not what she asked, and read strictly her question has no
answer at all. The mistake was in the answer, not in her question.

In the 1790s a new group appears. Miss Maria Middleton of Eden, near
Durham, answers at least four questions. Miss Betty Claxton of Benwell
gives the first printed answer to a problem about a window and a wall.
Miss Sarah Cowen sets two questions, and a Miss Castieau from Shropshire
answers one of them. In the 1800s Miss Mary Weston of Gainsborough
answers eight questions, more than any woman since the 1720s. In 1802 a
Mrs Macdonald of Spindlestone, near Bamburgh, asks how large a ball of
cork must be to fall at a steady ten feet a second through air: a
question about air resistance that would not look out of place in a
university mechanics course.

\hypertarget{nancy-mason}{%
\subsection{Nancy Mason}\label{nancy-mason}}

One of them we can see a little more clearly, thanks to the local
historian Sylvia Harrop.

Miss Nancy Mason of Clapham first appears in the Diary in 1792, aged
about nineteen. Her Clapham is the village under Ingleborough in the
Yorkshire Dales. She is in the Diary every year until 1798. In 1795 she
set a question whose answer was her own age on the first of January: 21
years, 11 months and 7 days. In 1796 she set one about a father who
leaves a square field of thirty acres to his five sons, to be divided
into strips of equal width running all round, one inside the other, the
eldest to have eight acres, the next seven, and so on down to four for
the youngest, whose share is the square in the middle. It is a lovely
problem, and it is in the fourth booklet.

She did not only do mathematics. She wrote verse charades, and long,
opinionated answers to the Diary's ``queries'': whether a virtuous old
maid or a young widow would make the more charming wife, where the
custom of valentines came from. In 1793 she shared a prize with a
clergyman, the Reverend Mr Ewbank, six free copies of the Diary each. In
1796 Ewbank printed a poem to her, offering, half in jest, to ``tie the
knot'' if he ever visited.

In 2010 Harrop went looking for her in the North Craven Heritage Trust
Journal. The parish registers of Clapham, Giggleswick and Bentham have
no baptism and no burial for any Nancy Mason. Nor do the Quaker records.
Harrop concludes that she was not from the district, probably came from
a comfortable or clerical family, and was staying with relatives or
friends. A modern reader judged her mathematics to be of about today's
A-level standard.

After 1798 she disappears. Harrop thinks she married. The Diary had
noticed the pattern itself: in 1792 it remarked on three ``fair and
ingenious sisters'' who had stopped writing one by one, as each of them
married. It may be the commonest ending in this whole story. We simply
cannot see it.

\hypertarget{a-warning-about-names}{%
\subsection{A warning about names}\label{a-warning-about-names}}

Not every woman's name was a woman. The most prolific ``woman'' setter
of the mid-century, Miss Ann Nicholls, who set ten questions in the
1760s, was in fact William Wales, the astronomer who later sailed with
Captain Cook. Leybourn knew him and says so in his index. So the early
pen-names, Xantippe, Adrastea, Silvia, may hide men too.

It cuts the other way as well. Leybourn collected only the mathematics.
The riddles, rebuses and verse, where women were far more numerous, are
not in his volumes. And his index lists only the solvers whose answers
he printed. The Diary's own long lists of those who ``also answered''
are left out. A surviving copy of the 1780 Diary names a Mrs Eleanor
Suggett as solving three of that year's questions. She is in none of
Leybourn's lists. The women in this story are a lower bound.

There is one more. When the editor Henry Beighton died in 1743, his
widow Elizabeth carried the Diary on, and the Stationers' Company
recognised her as its author for the next sixteen years, until her own
death in 1759. She edited the riddles herself, as she had done while her
husband was alive. The mathematics she left to others: first Anthony
Thacker, until his death in 1744, then Robert Heath, then Thomas
Simpson. The histories written afterwards gave the men's dates and left
her out. So for sixteen years a woman ran the Diary, in exactly the
years when women had stopped setting its questions.

\hypertarget{after-leybourn}{%
\subsection{After Leybourn}\label{after-leybourn}}

Leybourn's volumes stop in 1816, but the Diary ran on for another
twenty-four years, and bound copies of those issues survive in a few
libraries. The historians Joe Albree and Scott Brown counted only three
women who set mathematical questions between 1817 and 1836, and none
after. We think we have found them. Miss Mary Knight of Boston and Miss
Susan Jackson of Mile End each set a question in 1817. In 1834 a Miss L.
L. asked for the largest circle that fits in a triangle cut from a
semicircle, and two years later for the shortest route from London to
Constantinople that touches the equator.

The solvers are easier to find. Miss Sarah Ellis answered in 1824, and
went on answering after her marriage: the Diary for 1825 prints her as
``Mrs.~Sarah Clay (formerly Miss S. Ellis)'', and in 1828 Mrs S. Clay of
Moulton answered at least nine questions. In 1829 the first printed
answer to a question about nine numbers in arithmetic progression is by
Miss Jane Thoburn of Blyth, ``aged 12 years''. It is short, and it is
right.

By 1834 Miss L. L. was answering at least eight questions a year. The
editor, Olinthus Gregory of the Royal Military Academy, put a footnote
against her name: ``We admire the talents of our fair correspondent. If
she will only make up her mind to carry her own solutions through,
instead of referring to Simpson, Hymer, \&c.~she may soon vie with Donna
Agnesi, and our own Mrs.~Somerville.'' It is praise with a
schoolmaster's nudge in it. Her own answer to her circle question,
printed the next year, is exactly right.

In the last issue, for 1840, no woman appears in the mathematics at all.
The next year the Diary was merged with its rival, \emph{The Gentleman's
Diary}.

\hypertarget{what-we-made}{%
\subsection{What we made}\label{what-we-made}}

Working with Claude, I went through all four of Leybourn's volumes, one
surviving issue of the Diary for 1780, and three bound volumes of the
issues from 1816 to 1840, and produced nine booklets:

\begin{itemize}
\tightlist
\item
  \textbf{Book One:} 102 problems from 1707--1748, sorted by age from 5
  to 17, following the English school stages.
\item
  \textbf{Book Two:} 57 problems from 1802--1816.
\item
  \textbf{Book Three:} 55 problems from 1749--1776.
\item
  \textbf{Book Four:} 42 problems from 1777--1802.
\item
  \textbf{The Diary for 1780:} the ten questions of a single year, as
  one reader would have met them.
\item
  \textbf{After Leybourn:} 26 problems from 1816--1825.
\item
  \textbf{The Gregory years:} 22 problems from 1826--1834.
\item
  \textbf{The last years:} 20 problems from 1835--1840.
\item
  \textbf{Beyond A-level:} 55 harder problems from all four volumes,
  including the mechanics: pendulums, hanging chains, rolling balls, a
  falling Moon, a ball dropped through the centre of the Earth, and a
  spider chasing a fly that the Diary itself never managed to solve.
\end{itemize}

That is 389 problems, from every year the Diary ran. There are also
notes on the women of the Diary, with everything above and its sources.

As far as I know, this is the largest modern edition of the Diary's
mathematics, and the first collection of any kind of the questions set
after 1816, which Leybourn's volumes do not reach. Frank Swetz's history
of the Diary reprints a selection of its problems with solutions, and is
the place to go for the story of the almanac itself. These booklets are
meant to sit beside it: free, retold, and sorted for use in school.

Every problem is retold in modern English, with old money and measures
converted where it matters, and every answer has been recomputed. Along
the way we found that the Diary was sometimes wrong, and each booklet
has a small box saying so. One solver claimed a fares problem had
``innumerable answers''; it has exactly eight. Another printed a
pendulum length to twenty decimal places; the last four are wrong. Where
two solvers disagreed, we say who was right.

Claude drafted the solutions and retelling, recomputed every answer, ran
OCR on the scans and typeset everything in LaTeX. I chose the problems,
made the editorial decisions and checked the result. Each booklet says
so.

\hypertarget{why-it-matters}{%
\subsection{Why it matters}\label{why-it-matters}}

None of these women could have gone to a university, and few can have
been taught any mathematics at all. What they had was an annual almanac,
costing a few pence, that printed hard questions, printed the names of
the people who answered them, and corrected its mistakes in public the
following year.

That was enough. For a few decades it was enough for one question in
ten.

We spend a great deal of effort now on encouraging girls into
mathematics. We usually reach for the great names, for Ada Lovelace and
Mary Somerville and Emmy Noether. They are the right names. But there is
something to be said for Mary Wright and Nancy Mason as well. They were
not geniuses. They were not professionals. They were women who liked a
good problem, found a place that would print one, and sent their answers
in.

\hypertarget{where-to-find-it}{%
\subsection{Where to find it}\label{where-to-find-it}}

The LaTeX sources and PDFs are free on GitHub:
https://github.com/leznikm8/Ladies\_Diary. Use them in a classroom,
print them for your children, or fix my mistakes.

A magazine for women published advanced mathematics problems for 137
years. The least we can do is make them readable again.

\begin{center}\rule{0.5\linewidth}{0.5pt}\end{center}

\textbf{Sources.} T. Leybourn (ed.), \emph{The Mathematical Questions
Proposed in the Ladies' Diary} (4 vols, 1817). S. Harrop, ``The Enigma
of Miss Nancy Mason of Clapham'', \emph{North Craven Heritage Trust
Journal} (2010). S. A. Costa, \emph{The Ladies' Diary: Society, Gender
and Mathematics in England, 1704--1754} (PhD dissertation, Cornell,
2000). S. Costa, ``The Ladies' Diary: Gender, Mathematics, and Civil
Society in Early-Eighteenth-Century England'', \emph{Osiris} 17 (2002).
T. Perl, ``The Ladies' Diary or Woman's Almanack, 1704--1841'',
\emph{Historia Mathematica} 6 (1979). J. Albree and S. H. Brown, ``A
valuable monument of mathematical genius'', \emph{Historia Mathematica}
36 (2009). \emph{The Ladies' Diary} for 1816--1825 and 1825--1834
(digitised by Google Books) and for 1835--1840 (Harvard College Library
copy, digitised by HathiTrust). F. J. Swetz, \emph{The Impact and Legacy
of The Ladies' Diary} (MAA Press, 2021). S. E. Despeaux, ``Constance
Marks and the Educational Times'', in \emph{Women in Mathematics}
(Springer, 2017).

\emph{Correction, 30 October 2026: an earlier version of this article
said that Mary Wright answered the steeple question in 1709 and had
altered its figures herself. The question was set in 1710 and answered
in 1711, and it was the editor, John Tipper, who altered the figures. I
am grateful to Shelley Costa for the correction; see her dissertation,
pp.~124--131.}
\end{document}
